\section{Descomposici\'on espectral bicapa del operador angular}

Sea \(R=R^*=R^{-1}\) la involuci\'on activa y

\[
P_\pm=\frac{I\pm R}{2},
\qquad
T=q_+P_+ +q_-P_-,
\qquad 0<q_+<q_-<1.
\]

Con
\[
x=A_{\rm rad}=\frac{\pi A}{180},\qquad
y=C^*_{\rm rad}=\frac{\pi C^*}{180},
\]
los dos canales se construyen antes de evaluar el vac\'io:
\begin{equation}
q_+=e^{-x-y},\qquad q_-=e^{-x+y}.
\label{eq:vacuum-angular-channels}
\end{equation}
El intercambio de hojas invierte \(C^*\), es decir \(y\mapsto-y\), y permuta \(q_+\leftrightarrow q_-\). La pareja es as\'i la lectura espectral de un mismo antecedente angular, no dos par\'ametros ajustados por separado.

\begin{definition}[Respuesta constitutiva]
\begin{equation}
\mathcal V_{90,120}=(I-T^{90})(I-T^{120})^{-1}.
\label{eq:vacuum-operator}
\end{equation}
\end{definition}

Como \(\|T\|<1\), el denominador es positivo e invertible. El c\'alculo funcional da

\begin{equation}
\mathcal V_{90,120}=r_+P_+ +r_-P_-,
\qquad
r_\pm=\frac{1-q_\pm^{90}}{1-q_\pm^{120}}.
\label{eq:vacuum-eigenvalues}
\end{equation}

La afirmaci\'on de invertibilidad merece hacerse expl\'icita. Para cada hoja,
\(0<q_\pm^{120}<1\); por tanto
\[
I-T^{120}=(1-q_+^{120})P_++(1-q_-^{120})P_-
\]
es estrictamente positivo y
\[
(I-T^{120})^{-1}
=\frac{P_+}{1-q_+^{120}}+
 \frac{P_-}{1-q_-^{120}}.
\]
La respuesta \(\mathcal V_{90,120}\) no se obtiene resolviendo cuatro
ecuaciones escalares separadas: es el resultado del c\'alculo funcional de una
sola contracci\'on positiva. Esta observaci\'on impide que permeabilidad,
permitividad, impedancia y propagaci\'on pierdan el antecedente que comparten.

\begin{theorem}[Determinaci\'on espectral conjunta de las cuatro relaciones constitutivas]
Los cuatro caracteres constitutivos
\begin{equation}
\widehat\varepsilon=r_+^2,\qquad
\widehat\mu=r_-^2,\qquad
\widehat Z=\frac{r_-}{r_+},\qquad
\widehat c=\frac1{r_+r_-}
\label{eq:vacuum-four}
\end{equation}
son funciones espectrales de \(\mathcal V_{90,120}\). Satisfacen
\begin{equation}
\widehat\mu\widehat\varepsilon=\widehat c^{-2},
\qquad
\frac{\widehat\mu}{\widehat\varepsilon}=\widehat Z^2.
\label{eq:vacuum-relations}
\end{equation}
\end{theorem}

\begin{proof}
Las identidades se obtienen sustituyendo \cref{eq:vacuum-four}. No interviene ninguna carta dimensional ni ning\'un valor exterior.
\end{proof}

\begin{corollary}[Reconstrucci\'on de las dos hojas]
Los cuatro caracteres no contienen cuatro grados de libertad. Cualquiera de
los pares \((\widehat\mu,\widehat\varepsilon)\),
\((\widehat Z,\widehat c)\) o \((\mathcal E,\mathcal B)\) reconstruye los
dos autovalores positivos:
\begin{align}
r_-&=\sqrt{\widehat\mu}
=\sqrt{\widehat Z/\widehat c},
&r_+&=\sqrt{\widehat\varepsilon}
=\frac1{\sqrt{\widehat Z\widehat c}},
\label{eq:vacuum-reconstruct-r}\\
\log r_-&=\frac{\mathcal E+\mathcal B}{2},
&\log r_+&=\frac{\mathcal E-\mathcal B}{2}.
\label{eq:vacuum-reconstruct-log}
\end{align}
En particular, la orientaci\'on se pierde si se conserva s\'olo el producto,
pero se recupera al a\~nadir el cociente.
\end{corollary}

\begin{proof}
La primera l\'inea es la ra\'iz positiva de las identidades
\cref{eq:vacuum-four,eq:vacuum-relations}. La segunda invierte el sistema
lineal
\(\mathcal E=\log r_-+\log r_+\),
\(\mathcal B=\log r_--\log r_+\).
\end{proof}

Este corolario expresa una propiedad central de HMT: producto y cociente son
funcionales complementarios de un mismo estado. El producto determina la
propagaci\'on com\'un, mientras que el cociente conserva la orientaci\'on
relativa de las hojas interior y exterior.

En el certificado V2,

\begin{align*}
r_+&=0.9999996285482493631786952281\ldots,\\
r_-&=0.9997241371895183641315165853\ldots,\\
\widehat\varepsilon&=0.9999992570966367027604416155\ldots,\\
\widehat\mu&=0.9994483504793269350899815559\ldots,\\
\widehat Z&=0.9997245085389372154555543466\ldots,\\
\widehat c&=1.0002763104861575261063862689\ldots.
\end{align*}

Estos numerales son el reconocimiento terminal de \cref{eq:vacuum-operator}; no definen sus exponentes \(90,120\).

\section{Factorizaci\'on arm\'onica de los sectores \(90\) y \(120\)}

El semigrupo arm\'onico es

\[
\langle90,120\rangle=30\langle3,4\rangle.
\]

Si \(s=q^{30}\), entonces

\begin{equation}
\frac{1-q^{90}}{1-q^{120}}
=\frac{1-s^3}{1-s^4}
=\frac{1+s+s^2}{1+s+s^2+s^3}.
\label{eq:three-four-completion}
\end{equation}

La identidad tambi\'en separa la respuesta de su defecto. Si
\[
F_{3\to4}(s)=\frac{1+s+s^2}{1+s+s^2+s^3},
\qquad 0<s<1,
\]
entonces
\begin{equation}
d_{3\to4}(s):=1-F_{3\to4}(s)
=\frac{s^3}{1+s+s^2+s^3}>0.
\label{eq:vacuum-defect34}
\end{equation}
El t\'ermino adicional se conserva exactamente como defecto
positivo. Adem\'as,
\begin{equation}
F_{3\to4}'(s)
=-\frac{s^2(3+2s+s^2)}{(1+s+s^2+s^3)^2}<0.
\label{eq:vacuum-monotonic34}
\end{equation}
Por consiguiente, el orden de los canales angulares determina un\'ivocamente
el orden de las respuestas constitutivas. La asignaci\'on de
\(\widehat\mu\) y \(\widehat\varepsilon\) a las hojas respectivas se obtiene
anal\'iticamente, sin recurrir a sus expansiones decimales.

Cada hoja del vac\'io constituye la completaci\'on exacta de tres a cuatro
t\'erminos de un mismo modo elemental \(30\). En consecuencia, permeabilidad
y permitividad se interpretan como dos funcionales cuadr\'aticos de una misma
completaci\'on evaluada sobre hojas conjugadas, y no como par\'ametros
constitutivos independientes.

Como la funci\'on \((1-q^{90})/(1-q^{120})\) decrece en \(0<q<1\) y \(q_+<q_-\), se obtiene

\[
0<r_-<r_+<1,\qquad
0<\widehat\mu<\widehat\varepsilon<1,\qquad
0<\widehat Z<1<\widehat c.
\]

Los l\'imites del funcional determinan su geometr\'ia. Cuando \(q\to0^+\),
\(F(q)\to1\): las cuatro lecturas se acercan a la unidad y el defecto de
memoria desaparece. Cuando \(q\to1^-\),
\[
F(q)\longrightarrow\frac{90}{120}=\frac34.
\]
La respuesta toma, por tanto, valores en el intervalo positivo \((3/4,1)\). El
extremo \(3/4\) no es una constante ajustada: es el cociente de las dos
longitudes de canal y la realizaci\'on continua de la completaci\'on
tres--a--cuatro.

\section{Involuci\'on de hojas y conservaci\'on de la orientaci\'on}

El intercambio de hojas \(C^*\mapsto-C^*\) act\'ua por

\begin{equation}
\widehat\mu\longleftrightarrow\widehat\varepsilon,\qquad
\widehat Z\longmapsto\widehat Z^{-1},
\qquad
\widehat c\longmapsto\widehat c.
\label{eq:vacuum-sheet}
\end{equation}

La propagaci\'on com\'un es par; la impedancia conserva la orientaci\'on.
Esta distinci\'on reaparecer\'a en Barbero y CKM: no todos los funcionales
asociados al cambio de hoja tienen la misma paridad.

\section{Caracteres logar\'itmicos y compatibilidad tracial}

Definamos

\[
\mathcal E=\log(r_+r_-),\qquad
\mathcal B=\log(r_-/r_+).
\]

Con la traza normalizada \(\tau=\Tr/2\),

\begin{equation}
2\tau(\log\mathcal V)=\mathcal E,
\qquad
-2\tau(R\log\mathcal V)=\mathcal B.
\label{eq:vacuum-traces}
\end{equation}

La pareja \((\mathcal E,\mathcal B)\) constituye un sistema de coordenadas lineal del logaritmo
operatorio:
\begin{equation}
\log\mathcal V
=\frac{\mathcal E}{2}I-\frac{\mathcal B}{2}R.
\label{eq:vacuum-log-decomposition}
\end{equation}
La componente escalar es invariante bajo el intercambio de hojas; la
componente orientada cambia de signo. Esta descomposici\'on demuestra
directamente la paridad de la propagaci\'on y el car\'acter orientado de la
impedancia.

La amplificaci\'on non\'adica

\[
\mathcal V_m=\mathcal V\otimes I_{9^m}
\]

es compatible con las expectativas \(E_m=\mathrm{id}\otimes\tau_9\):

\[
E_m(\mathcal V_{m+1})=\mathcal V_m.
\]

Por tanto, \(\mathcal E\) y \(\mathcal B\) no son peculiaridades de una matriz \(2\times2\); permanecen en la torre UHF y en su cierre tracial \(II_1\).

M\'as precisamente, para todo polinomio \(p\) y despu\'es, por continuidad,
para toda funci\'on continua en el espectro,
\begin{equation}
E_m\bigl(p(\mathcal V_{m+1})\bigr)=p(\mathcal V_m).
\label{eq:vacuum-functional-refinement}
\end{equation}
La compatibilidad no preserva solamente dos n\'umeros: preserva todo el
\'algebra funcional generada por la respuesta. En particular conserva sus
proyectores espectrales, sus potencias, su logaritmo y las formas cuadr\'aticas
que miden los defectos \cref{eq:vacuum-defect34}.

\section{Realizaci\'on dimensional de las relaciones constitutivas}

La Mec\'anica Dimensional determina

\begin{align}
Z_{\rm vac}&=\widehat Z\,u_Su_Q^{-2},\\
c_{\rm vac}&=\widehat c\,u_C,\\
\mu_{\rm vac}&=\widehat\mu\,u_Su_C^{-1}u_Q^{-2},\\
\varepsilon_{\rm vac}&=\widehat\varepsilon\,u_S^{-1}u_C^{-1}u_Q^2.
\label{eq:vacuum-sections}
\end{align}

Las identidades constitutivas se mantienen en las rectas:

\begin{equation}
\mu_{\rm vac}\varepsilon_{\rm vac}=c_{\rm vac}^{-2},
\qquad
\mu_{\rm vac}/\varepsilon_{\rm vac}=Z_{\rm vac}^2.
\end{equation}

La carta SI es posterior. No es necesario declarar \(c\) primitivo y resolver despu\'es \(\mu\varepsilon=c^{-2}\): los tres objetos descienden conjuntamente de \(\mathcal V\) y de las rectas declaradas.

La separaci\'on de niveles se expresa mediante el diagrama
\[
\begin{array}{c}
(A,C^*)\longrightarrow T\longrightarrow\mathcal V
\longrightarrow
(\widehat\mu,\widehat\varepsilon,\widehat Z,\widehat c)
\\[2mm]
\Big\downarrow\text{ cartas de recta}\hspace{36mm}
\Big\downarrow\text{ evaluaci\'on final}
\\[2mm]
(u_S,u_Q,u_C)\longrightarrow
(\mu_{\rm vac},\varepsilon_{\rm vac},Z_{\rm vac},c_{\rm vac})
\longrightarrow\text{ coordenadas SI}.
\end{array}
\]
La primera l\'inea es adimensional y operatoria; la segunda determina tipos de
magnitud; la tercera selecciona una carta metrol\'ogica. Alterar la carta final
no modifica ni los autovalores ni las identidades constitutivas. En ese
sentido preciso, la unificaci\'on matem\'atica--f\'isica consiste en distinguir
el invariante geneal\'ogico de su sistema de coordenadas, y no en suprimir la
estructura dimensional.

\begin{proposition}[Minimalidad espectral]
Sea \(W\) un operador positivo de dos hojas con espectro simple
\(\{r_+,r_-\}\). Si cuatro caracteres positivos satisfacen
\(\mu\varepsilon=c^{-2}\) y \(\mu/\varepsilon=Z^2\), si la propagaci\'on
com\'un se determina mediante \(c^{-1}=\det W=r_-r_+\), y si su orientaci\'on
fija \(Z=r_-/r_+\), entonces necesariamente
\[
(\mu,\varepsilon,Z,c)
=\left(r_-^2,r_+^2,\frac{r_-}{r_+},\frac1{r_-r_+}\right).
\]
Por tanto \cref{eq:vacuum-four} no es una parametrizaci\'on entre muchas:
es la realizaci\'on positiva m\'inima compatible con producto, cociente y
orientaci\'on.
\end{proposition}

\begin{proof}
De \(Z=r_-/r_+\) y \(c^{-1}=\det W=r_-r_+\) se obtienen de forma \`unica las
ra\'ices positivas. Las dos identidades constitutivas fuerzan despu\'es
\(\mu=r_-^2\) y \(\varepsilon=r_+^2\).
\end{proof}

